% Figure: parameter-identifiability map. Three regions in (r, K) space
% for the logistic equation, showing how the data window selects which
% parameters are tightly identified.
\begin{figure}[htbp]
\centering
\begin{tikzpicture}[font=\small]
% Axes
\draw[->, thick, draw=engblack!70] (0, 0) -- (8.5, 0) node[right, font=\small] {$r$ (intrinsic growth rate, $\si{\per\hour}$)};
\draw[->, thick, draw=engblack!70] (0, 0) -- (0, 5.5) node[above, font=\small, align=center] {$K$\\ (carrying capacity)};
% Tick labels
\foreach \xt/\xv in {1.5/0.5, 3.0/1.0, 4.5/1.5, 6.0/2.0, 7.5/2.5} {
  \draw[draw=engblack!70] (\xt, 0) -- (\xt, -0.08);
  \node[font=\scriptsize, anchor=north] at (\xt, -0.1) {\xv};
}
\foreach \yt/\yv in {1.0/1, 2.0/5, 3.0/10, 4.0/50} {
  \draw[draw=engblack!70] (0, \yt) -- (-0.08, \yt);
  \node[font=\scriptsize, anchor=east] at (-0.1, \yt) {\yv};
}

% Region A: data span only exponential regime; K barely identified
\fill[engred!15, opacity=0.7] (0.4, 0.4) rectangle (8.0, 4.8);
\draw[draw=engred!60, dashed, thick] (0.4, 0.4) rectangle (8.0, 4.8);

% Region B: data span the inflection; both identified
\fill[engblue!20, opacity=0.7] (2.0, 1.5) rectangle (5.5, 3.5);
\draw[draw=engblue!70, thick] (2.0, 1.5) rectangle (5.5, 3.5);

% Region C: data span saturation only; r barely identified
\fill[englime!25, opacity=0.7] (3.0, 2.0) ellipse (1.5 and 0.8);
\draw[draw=engblack!70, thick] (3.0, 2.0) ellipse (1.5 and 0.8);

% Labels
\node[font=\footnotesize, align=center, color=engblack] at (6.5, 4.2)
  {(A) data span\\ exponential only:\\ $r$ identified, $K$ wide};
\node[font=\footnotesize, align=center, color=engblack] at (3.75, 2.5)
  {(B) data span\\ inflection: both identified};
\node[font=\footnotesize, align=center, color=engblack, fill=white, inner sep=1pt] at (3.0, 0.9)
  {(C) data span saturation: $K$ identified, $r$ wide};
\end{tikzpicture}
\caption[Parameter identifiability and the data window]{The
identifiability of $(r, K)$ in the logistic model depends on the
window of time over which the data are collected. (A) Data confined to
the early exponential phase identify $r$ tightly (it sets the early
slope on a log scale) but leave $K$ almost unconstrained (the
saturation has not been observed). (B) Data spanning the inflection
point (population near $K/2$) identify both: the inflection's location
fixes $K$ and the slope through it fixes $r$. (C) Data confined to
the saturation tail identify $K$ tightly but leave $r$ poorly
constrained because the transient has already passed. The identifiability
diagnostic is therefore part of the experimental design: choose the
data window such that the parameters the engineer needs are in
region~(B). The discipline appears in Cobelli and Distefano's
identifiability literature \cite{paper:v2ch18-cobelli-distefano-1980}
and underwrites Exercise 18.18 (design of an identification experiment).}
\label{fig:vol02:ch18:identifiability}
\end{figure}
